\section{Conclusion}
\label{sec:conclusion}

\Profile{} combines a narrowly defined private-state relation, a
circle-domain transparent argument, and a Solana account transition whose
proof verification and public-state mutation occur atomically.  For the
frozen (q=18,g=37) instance, the conservative classical-ROM analysis gives
a 100.161-bit work-normalized argument-soundness floor under the stated
finite-parameter and state-restoration premises.  Under the affine-image
premise, the declared programmable-ROM Proof-or-Abort view has a witness-free
simulator with a 104.024-bit real-versus-simulator floor and a 103.024-bit
pairwise-witness floor.

The released \ProfileProofBytes-byte proof verifies under a
\ProfileSbfBytes-byte SBF program.  The maximum measured local tag~65 path is
\ProfileComputeUnits{} CU.  The exact release subsequently executed on
mainnet-beta: transaction
\nolinkurl{3G1voggszvDMGi5PbGM1kuEMYKvh2TNMbH6hHHwndUdRQJNT7ehRFpQpksxLnx5tp2xkS5jGi359rVXk42sRPFcv}
consumed \ProfileMainnetComputeUnits{} CU, advanced the pool sequence,
created the canonical nullifier marker, and refunded the finalized proof
account.  Independent finalized RPC records agree on the transaction and
post-state.  Program deployment, proof transport, and finalization were
separate setup operations; ProgramData recovery and the treasury sweep were
separate cleanup operations.

The accompanying Lean~4 project kernel-checks the relation core, finite
security arithmetic, circle and masking lemmas, and capstone composition. The
current repository extends that approach with an Aeneas-generated
Rust-to-Lean path for the later V5 Tag-67 candidate, while keeping the V5
candidate distinct from the q18/g37 transaction reported here.

The result establishes feasibility for this one-input, one-output same-path
transition under the measured Solana limit.  Broader transaction relations,
lower prover latency, independent security review, and sustained deployment
engineering remain separate work.
